\documentclass[11pt]{article}
\usepackage{icmlko}
\begin{document}

\title{몇 개의 축이 갈리나 --- 쌍에서는 되고 도메인에서는 안 되며, 신뢰도로 쓰려니 정작 접전에서 평평하다}
\author{월드모델 실험실}
\date{노트 531 · 2026년 8월 4일}
\maketitle

\begin{abstract}
노트 530 의 분해능(구별 불가능 쌍)은 \emph{극단만} 세서 태거 몫을 $0.3$ \%p
하한으로만 줬다. 더 민감한 지표를 세웠다 --- \textbf{쌍마다 값 또는 결측
무늬가 다른 축의 수 $D$}.

\begin{center}
\begin{tabular}{lrrr}
\hline
축 차이 $D$ & 전체 & \textbf{접전} & 명백 \\
\hline
$1\!\sim\!4$ & $56.6\%$ & $53.3\%$ & $73.6\%$ \\
$5\!\sim\!7$ & $63.6\%$ & $55.7\%$ & $75.9\%$ \\
$8\!\sim\!11$ & $67.6\%$ & $56.9\%$ & $82.0\%$ \\
$\ge 12$ & $\mathbf{70.1\%}$ & $\mathbf{57.2\%}$ & $\mathbf{85.0\%}$ \\
\hline
폭 & $+13.5$ \%p & $\mathbf{+3.8}$ \textbf{\%p} & $+11.4$ \%p \\
\hline
\end{tabular}
\end{center}

사전 등록 ①(\emph{$D$ 클수록 정확}) \textbf{통과}, ②(\emph{분해능이 못 보는
곳도 본다}) \textbf{반증}. 그리고 덤으로 얻은 ``예측마다 붙는 신뢰도''는
\textbf{정작 접전에서 폭이 $3.8$ \%p 다.}
\end{abstract}

\begin{figure}[t]
\centering
\includegraphics[width=\textwidth]{figs/f1.pdf}
\caption{\textbf{가} --- 세 층에서의 $D$. \textbf{나} --- 쌍 수준과 도메인
수준. \textbf{다} --- 읽기.}
\end{figure}

\section{먼저, 내가 내 표를 잘못 읽을 뻔했다}

첫 실행의 구간 표는 $D=1$ 에서 $24.9\%$, $D \ge 12$ 에서 $45.2\%$ 를 냈다.
\textbf{전부 $50\%$ 아래}인데 전체 가중 정확도는 $67.2\%$ 다 --- 앞뒤가 안 맞는다.

\begin{quote}
가중 평균을 낼 때 \textbf{그 구간에 쌍이 없는 도메인의 무게를 안 뺐다.}
$D=1$ 쌍은 다섯 도메인에만 있는데 열한 도메인의 무게로 나눴으니 값이 눌린다.
고치니 $54.3\% \to 70.1\%$ 로 \textbf{단조}가 됐고 사전 등록 ①은 통과다.
\end{quote}

\textbf{$50\%$ 아래라는 것이 신호였다.} 그대로 ``반증''으로 적었으면 참인
명제를 죽일 뻔했다. 도메인 안에서도 $9/11$ 이 단조다(예외는 팝업 $\cdot$ 도서
로, 둘 다 유보가 작다).

\section{사전 등록 ② --- 반증}

도메인 수준으로 올리면 새 정보가 없다.

\begin{center}
\begin{tabular}{lr}
\hline
 & 순위상관 \\
\hline
중앙 $D$ ↔ 정확도 (전체 $11$) & $+0.572$ \\
중앙 $D$ ↔ 정확도 (분해능 $0.9$ 이상 여덟) & $\mathbf{-0.140}$ \\
\textbf{중앙 $D$ ↔ 분해능} & $\mathbf{+0.928}$ \\
\hline
\end{tabular}
\end{center}

사전 등록은 ``$0.9$ 이상 여덟에서 $\ge +0.30$''이었고 $-0.140$ 이 나왔다
(노트 530 의 분해능은 거기서 $+0.062$ 였다).

\begin{quote}
\textbf{두 지표가 같은 것을 재고 있다}($\rho = +0.928$). \textbf{태거 품질
지표는 문턱 지표를 못 벗어난다} --- 심한 곳(시장팝업 $\cdot$ 만화)을 고르는
데는 되고, 나머지의 \emph{정도}를 재는 데는 안 된다.
\end{quote}

\section{덤, 그리고 그 덤도 반쪽이다}

$D$ 는 \textbf{라벨을 안 보고} 계산된다. 그러니 예측마다 붙일 수 있는
\emph{신뢰도}가 된다 --- ``이 둘은 축 하나로만 갈립니다''와 ``열두 개가
갈립니다''는 다른 답이어야 한다. 전체 쌍에서 폭이 $+13.5$ \%p 니 쓸 만하다.

\begin{quote}
\textbf{그런데 접전에서 폭이 $\mathbf{+3.8}$ \textbf{\%p} 다}($53.3 \to 57.2$).
노트 528 이 보인 대로 \emph{실무의 판은 접전}인데, 거기서 $D$ 는
\textbf{불러도 되는 접전과 안 되는 접전을 못 가린다.} 신뢰도가 필요 없는
곳에서만 신뢰도가 작동한다.
\end{quote}

\section*{조항}

\begin{quote}
\textbf{지표가 쌍 수준에서 되는 것과 도메인 수준에서 되는 것은 다르다.}
쌍 안에서 강하게 단조인 양($54.3 \to 70.1$)이 도메인 사이에서는 기존 지표와
$\rho = +0.928$ 로 겹칠 수 있다. \textbf{어느 수준에서 쓸 것인지를 먼저 적고
그 수준에서 판정한다.}
\end{quote}

\section*{남은 것}

\begin{itemize}
\item \textbf{접전을 가르는 신호는 축 개수가 아니다.} 무엇인지 아직 모른다 ---
모델의 예측 분산(자루 $32$ 개의 흩어짐)이 후보다. 라벨을 안 보므로 같은 자리에
쓸 수 있다.
\item \textbf{시장팝업은 유보 쌍의 $83.6\%$ 가 $D \le 1$ 이다}(다음이 만화
$8.7\%$). 노트 530 의 진단이 다른 셈법에서도 같은 곳을 가리킨다.
\item \textbf{학습 쪽 $D$} 는 안 쟀다.
\end{itemize}

\end{document}
